\section{The canonical M\"obius geometry}
\label{sec:canonical-cloud}
%======================================================================

The M\"obius decomposition needs one reproducible point per squarefree source.  Choosing the largest prime factor supplies that point and gives a signed height that can be followed across cutoffs.

The smooth--transport decomposition does not use every point on each factor fibre.  It selects one canonical point per squarefree source.

For squarefree $m>1$, define
\begin{equation}
 q_m:=\Pplus(m),
 \qquad
 c_m:=\frac{m}{q_m}.
 \label{eq:canonical-cq}
\end{equation}
Unlike the ordered factor pair in \cref{sec:full-fibre}, the canonical pair is not reordered.  The cofactor $c_m$ may exceed $q_m$.  For example,
\[
 m=30,
 \qquad
 q_m=5,
 \qquad
 c_m=6.
\]
Therefore define the signed Fermat coordinate
\begin{equation}
 a_m:=\frac{c_m+q_m}{2},
 \qquad
 b_m:=\frac{q_m-c_m}{2}\in\frac12\Z,
 \label{eq:signed-Fermat}
\end{equation}
and the canonical squared point
\begin{equation}
 w_m:=(a_m+ib_m)^2
 =m+iY_m,
 \label{eq:canonical-w}
\end{equation}
where
\begin{equation}
 \boxed{Y_m=2a_mb_m=\frac{q_m^2-c_m^2}{2}.}
 \label{eq:signed-Y}
\end{equation}
The sign of $Y_m$ records the order of the canonical factors:
\begin{equation}
 Y_m>0\iff c_m<q_m,
 \qquad
 Y_m<0\iff c_m>q_m.
 \label{eq:Y-sign}
\end{equation}

\subsection{Parameterization by cofactors and primes}

Every squarefree $m>1$ has a unique representation
\begin{equation}
 m=cq,
 \label{eq:cq-general}
\end{equation}
where
\begin{equation}
 \mu(c)\ne0,
 \qquad
 q\text{ is prime},
 \qquad
 q>\Pplus(c).
 \label{eq:canonical-admissibility}
\end{equation}
Conversely, every pair satisfying \cref{eq:canonical-admissibility} produces a squarefree integer $m=cq$ with largest prime factor $q$.  Moreover,
\begin{equation}
 \mu(cq)=-\mu(c).
 \label{eq:canonical-sign}
\end{equation}

Define the canonical index set
\begin{equation}
 \mathscr C
 :=
 \left\{
 (c,q):
 \mu(c)\ne0,
 \ q\text{ is prime},
 \ q>\Pplus(c)
 \right\}.
 \label{eq:canonical-index}
\end{equation}
The squared point cloud is
\begin{equation}
 \mathscr W
 :=
 \left\{
 cq+i\frac{q^2-c^2}{2}:(c,q)\in\mathscr C
 \right\}.
 \label{eq:canonical-cloud}
\end{equation}
The signed point measure is
\begin{equation}
 \nu
 :=
 \delta_1
 -
 \sum_{(c,q)\in\mathscr C}
 \mu(c)\,
 \delta_{\,cq+i(q^2-c^2)/2}.
 \label{eq:signed-measure}
\end{equation}
The atom at $1$ records $\mu(1)=1$.

\begin{example}[A point below the real axis]
For $m=30$, the canonical pair is $(c,q)=(6,5)$, so
\[
 a=\frac{11}{2},
 \qquad
 b=-\frac12,
 \qquad
 w_{30}=30-\frac{11}{2}i.
\]
This point enters a square prefix only after the cutoff already exceeds $q=5$, so it is smooth immediately and never occupies the transport channel.
\end{example}

%======================================================================
